% concatenated from Armflds.tex
\documentclass[12pt,a4paper]{article}

%% The amssymb package provides various useful mathematical symbols
\usepackage{mathtools,amssymb,mathrsfs}
%% The amsthm package provides extended theorem environments
\usepackage{amsthm}
\usepackage{enumerate}
\usepackage{physics}
\usepackage[hidelinks]{hyperref}
\usepackage{algorithm}
\usepackage{algpseudocode}
\usepackage{graphicx}
\usepackage[backend=biber, style=numeric, sorting=nty]{biblatex}
%\usepackage{cleveref}
\usepackage{float}


\hypersetup{
        colorlinks   = true,
        citecolor    = blue,
        linkcolor=blue  %%%ADDED for the links (needed in e.g. article)
}

\newcommand{\footremember}[2]{%
    \footnote{#2}
    \newcounter{#1}
    \setcounter{#1}{\value{footnote}}%
}
\newcommand{\footrecall}[1]{%
    \footnotemark[\value{#1}]%
} 

\newtheorem{theorem}{Theorem}[section]
\newtheorem{lemma}[theorem]{Lemma}
\newtheorem{prop}[theorem]{Proposition}
\newtheorem{cor}[theorem]{Corollary}
\newtheorem{conjecture}[theorem]{Conjecture}
\theoremstyle{definition}
\newtheorem{definition}[theorem]{Definition}
\newtheorem{example}[theorem]{Example}
\theoremstyle{remark}
\newtheorem*{remark}{Remark}

\DeclareMathOperator{\Var}{Var}
\DeclareMathOperator{\Cov}{Cov}
\DeclareMathOperator*{\argmin}{argmin}
\DeclareMathOperator*{\arginf}{arginf}
\DeclareMathOperator{\dist}{dist}
\DeclareMathOperator{\Vol}{Vol}
\DeclareMathOperator{\dVol}{dVol}
\DeclareMathOperator{\id}{id}
\DeclareMathOperator{\St}{St}
\DeclareMathOperator{\Gr}{Gr}
\DeclareMathOperator{\inv}{inv}
\newcommand{\ddf}[3]{\frac{\dd^{#1} #2}{\dd #3^{#1}}}

\DeclareNameAlias{author}{last-first}

%\def\appendixname{\hspace{-5pt}}

\addbibresource{references.bib}

\title{Autoregressive Processes on Stiefel and Grassmann Manifolds}

\author{Jordi-Llu\'{i}s Figueras\footremember{upps}{Department of Mathematics, Uppsala University}\footremember{huawei}{Partially funded by Huawei technologies.} \\
\small{\texttt{jordi-lluis.figueras@math.uu.se}}
\and
Aron Persson\footrecall{upps} \footrecall{huawei} \footnote{Corresponding Author}   \\
\small{\texttt{aron.persson@math.uu.se}}
}





\begin{document}









\maketitle


\begin{abstract}
%% Text of abstract
System identification of autoregressive processes on Stiefel and Grassmann manifolds are presented and studied. We define the system parameters as elements in the orthogonal group and we show that the system can be estimated by averaging over observations. Then we propose an algorithm on how to compute these system parameters using conjugate gradient descent on Stiefel and Grassmann manifolds, respectively.
\end{abstract}


\textbf{Keywords} Statistics on manifolds; differential geometry; Homogeneous spaces; directional statistics

\textbf{2020 Mathematics subject classification} 62R30; 53Z50




%% \linenumbers

%% main text

\section{Introduction}
\label{sec:intro}

In modern signal processing and MIMO systems, a state is described by a general $n$ by $k$ matrix. But it is only the orthogonal part of the state that is measured, i.e. one measures the orthogonal part of the polar decomposition of a $n$ by $k$ matrix. In other words the measurements are points on Grassmann or Stiefel manifolds, depending on application and context, see for example \cite{Kim2020stochcons,kong2022momentum,li2013riemannian,schwarz2014predictive,solo2014approach,xue2020blind}.

Now consider a sequence of measurements over time, a time-series that takes values on a Grassmann or Stiefel manifold. If one now wants to make a prediction of future states or analyze the dynamics of such a time series, we have to make a suitable model for the time series. In $\mathbb{R}$-valued statistics, the most basic model used in time series analysis are the autoregressive processes. This then begs the question: What is a natural definition of an autoregressive process on Stiefel and Grassmann manifolds? Moreover, how can one estimate the parameters of the model in a reasonable manner?

In \cite{xavier2006generalization}, Xavier and Manton introduced some AR-processes on Riemannian manifolds, a treatise we draw inspiration from and aim to specialize to the class of homogeneous spaces. The goal is to enable us to estimate the system parameters on Grassmann and Stiefel manifolds. Since, by definition, a homogeneous space has a transitive Lie-group action, this naturally guides us to consider system parameters as elements in the Lie group. In Proposition \ref{prop:Argroup}, we show that one may model an AR($N$)-processes on Stiefel and Grassmann manifolds by considering the system parameter as an element in the orthogonal group, see Equation \eqref{eq:liegroupactionAR}. Moreover, by Theorem \ref{thm:ARexpect} we show that the system parameter can be estimated, and numerically computed, by averaging out over the observations. Using the definition for the expectation on manifolds, we propose a conjugate gradient descent method in Algorithm \ref{alg:stiefelgrad}, which utilize an approximation of the distance function for the Grassmann and Stiefel manifolds.

This paper is relevant to anyone doing directional statistics, the most common form is that of directional statistics on the spheres $\mathbb{S}^n$. The spheres are special cases of the Stiefel manifolds, more precisely $\mathbb{S}^n \cong \St_{n+1,1}$. Ever since Fischer in \cite{fisher1953dispersion} popularized the idea of modeling data on spheres in order to ascertain geological information on rocks, see also \cite{mardia2025fisher}, there has been an abundance of applications of directional statistics. Directional statistics has found its way into meteorology \cite{bowers2000directional,li2018wavefront,nunez2015bayesian,Nunez-Gutierrez-Escarela-2011}, into medicine \cite{demir2019application,karaibrahimoglu2021circular} and into machine learning \cite{kanamori2012non}, to name a few.


We start, in Section \ref{sec:AutRM}, by giving an overview on what has been done previously in the area of autoregressive processes on Riemannian manifolds, and general statistics on manifolds. Then, in Section \ref{sec:ARhomo}, we specialize to homogeneous spaces and we give some new results. Thereafter, in Section \ref{sec:conjgrad}, we give some algorithms that estimate and approximate the system parameters on the Stiefel and Grassmann manifold using the conjugate gradient method. A fundamental introduction to Riemannian geometry is given in Appendix \ref{sec:basicriemgeom} that introduces concepts needed in order to understand the reasoning in this paper. In the second part of the appendix, Appendix \ref{sec:simulations}, one can find some figures showing simulations for a select few Stiefel and Grassmann manifolds illustrating the algorithms given in Section \ref{sec:conjgrad}.


\section{Statistics on Riemannian Manifolds}
\label{sec:AutRM}

Throughout, let $(\Omega, \mathcal{A},\mathbb{P})$ be a probability space. Let $M$ be a Riemannian manifold. We shall recall a slightly modified, but still equivalent, version of the "wrapped process" definition of an AR-process found in \cite{xavier2006generalization}. Denote $\exp_p$ and $\log_p$ to be the Riemannian exponential map centered at $p$ and its inverse, respectively, see Definition \ref{def:geodexplog}. Moreover, the parallel transport from $p\in M$ to $q\in M$ along a geodesic from $p$ to $q$ is written as $P_{p,q}:T_pM\rightarrow T_q M$.

\begin{definition}
\label{def:arprocess}
Let $M$ be a connected and geodesically complete Riemannian manifold, and let $Z_\ell\in M$, $\ell\in \mathbb{Z}$ be a time-series. Let $N \in \mathbb{N}$, and suppose there are $N$ (1,1)-tensor fields $A_j:TM\rightarrow TM$, i.e. for every $p\in M$ it holds that
\[
A_j\mid_{T_p M}:T_pM \rightarrow T_pM
\]
is linear. Then, an \emph{AR($N$)-process} on $M$ is defined as
\[
Z_\ell = \exp_{Z_{\ell-1}}\left(\varepsilon_\ell+ \sum_{j=1}^N A_j P_{Z_{\ell-j},Z_{\ell-1}}( \log_{Z_{\ell-j}}(Z_{\ell-j-1})) \right),
\]
where $\varepsilon_\ell:\Omega \rightarrow T_{Z_{\ell-1}}M$ is a random variable with $\mathbb{E}[\varepsilon_\ell]=0$, $\Cov(\varepsilon_\ell)=\id_{T_{Z_{\ell-1}}M}$.
\end{definition}

For our purposes we need the notion of expected value of a Riemann manifold-valued random variable. This is precisely the well known least square minimizer on $\mathbb{R}^n$ adapted and generalized to Riemannian manifolds. Suppose $\dist$ is the Riemannian distance function on $M$ and $\Vol$ is the Riemannian volume measure on $M$. Following \cite{dubey2019frechet,pennec2006intrinsic}, the expectation of an absolutely continuous random variable $X:(\Omega,\mathcal{F},\mathbb{P}) \rightarrow M$ is then any point $y\in M$ satisfying
\begin{equation}
\mathbb{E}[X] = \arginf_{y\in M} \int_M \dist^2(y,x)p_X(x) \dVol(x).
\label{eq:expectation}
\end{equation}
Here $p_X = \frac{\dd X_*\mathbb{P}}{\dVol}$ is the Radon-Nikodyn derivative of the push forward measure, defined by $X_*\mathbb{P}(A)= \mathbb{P}(X^{-1}(A))$ for any Borel subset $A\subseteq M$, with respect to the Riemannian volume measure. Note that in contrast to random variables on $\mathbb{R}^n$, the expected value need not be uniquely defined. This is clearly illustrated by considering a uniformly distributed random variable on $\mathbb{S}^2$, i.e. constant probability density function, then any, and thus all, points on $\mathbb{S}^2$ minimize the integral above.

\section{AR-Processes on Homogeneous Spaces by Lie Group Action}
\label{sec:ARhomo}

Let $G$ be a compact Lie group with bi-invariant Riemannian metric, and let $M$ be a geodesically complete $G$-homogeneous space with metric induced by $G$. See Appendix \ref{sec:basicriemgeom} for the definition.

\begin{prop}
\label{prop:Argroup}
Suppose $Z_\ell$ is an AR($N$)-process on $M$ as in Definition \ref{def:arprocess} with system parameters $A_j$.
Then, for every $1\leq j\leq N$ and any pair of fixed points $Z_{\ell-j},Z_{\ell-j-1}\in M$ there is a $\Phi_j\in G$ such that
\begin{equation}
\Phi_j Z_{\ell-j} = \exp_{Z_{\ell-j}}(A_j \log_{Z_{\ell-j}}(Z_{\ell-j-1})).
\label{eq:liegrouptensorfield}
\end{equation}
\end{prop}

\begin{proof}
Denote
\[
q = \exp_{Z_{\ell-j}}(A_j \log_{Z_{\ell-j}}(Z_{\ell-j-1})),
\]
and
\[
\pi: G\rightarrow M=G/H
\]
as the quotient projection from $G$ to $M$. Note that since $\pi$ is defining the Lie-group action on $M$, we have 
\[
\Phi_j Z_{\ell-j} = \Phi_j \pi(\Psi)=\pi(\Phi_j \Psi)
\]
where $\Phi_j,\Psi\in G$ and $\Psi$ is a representative for $Z_{\ell-j}$. Suppose $\xi\in G$ is such that $\pi(\xi)=q$ then taking $\Phi_j = \Psi \xi^{-1}$ is a Lie group element such that the equality in \eqref{eq:liegrouptensorfield} holds true.
\end{proof}

Note that the above may also be concluded by observing that the Lie group action on homogeneous spaces is transitive.

\begin{remark}
One important reason one may want to consider system parameters as Lie group elements is that for Stiefel and Grassmann manifolds the degrees of freedom for finding parameters $\Phi \in O(n)$ is much smaller than the degrees of freedom for finding tensor fields $A_j\in \mathcal{T}^1_1(M)$ so they will in general be less taxing to compute. Moreover, what the proof above indicates is that due to a high degree of multicollinearity of the system parameters, we can without loss of generality assume that the system parameters are elements in the Lie group instead.
\end{remark}

Therefore, we henceforth consider autoregressive processes of the following form
\begin{equation}
Z_\ell = \exp_{Z_{\ell-1}}\left(\varepsilon_\ell + \sum_{j=1}^N P_{Z_{\ell-j},Z_{\ell-1}}\log_{Z_{\ell-j}}(\Phi_j Z_{\ell-j}) \right)
\label{eq:autoregressivehomo}
\end{equation}
and with $\Cov(\varepsilon_\ell)=\sigma^2 \id_{T_{Z_{\ell-1}}M}$.

In what follows, we want to ensure that for an AR(1)-processes with reasonable assumptions on the noise, on average satisfy $Z_\ell=\Phi Z_{\ell-1}$. We shall impose a relatively weak symmetry condition on the noise, the condition of \textit{reflection symmetry}.

\begin{definition}
Let $M$ be a complete Riemannian manifold. A function $f:M\rightarrow M$ is said to be \textit{reflection symmetric} around a point $x\in  M$ if for all $V\in T_x M$ it holds that
\begin{equation}
\dist(f(\exp_{x}(V)),f(x) )=\dist(f(\exp_{x}(-V)),f(x)).
\label{eq:liegroupactionAR}
\end{equation}
Similarly a function $s:M \rightarrow \mathbb{R}$ is said to be \textit{reflection symmetric} around $x\in M$ if
\[
s(\exp_x(V))=s(\exp_x(-V))
\]
for all $V\in T_x M$.
\end{definition}

Next is a fairly general lemma which connects reflection symmetric functions $f:M\rightarrow M$ and reflection symmetric probability density functions.

\begin{lemma}
\label{lem:lipreflect}
Let $M$ be a Riemannian manifold, and let
\[
X:(\Omega,\mathcal{F},\mathbb{P}) \rightarrow M
\]
be an absolutely continuous random variable. Suppose $m=\mathbb{E}[X]$ exists and is unique. Assume $p_X(y)$ is reflection symmetric around $m$. Let $f:M\rightarrow M$ be an injective, non-exhaustive and reflection symmetric function and assume its inverse $f^{-1}:M\rightarrow M$ is reflection symmetric as well.
Then, the probability density function of $f(X)$ is reflection symmetric around $f(m)$.
\end{lemma}

\begin{proof}
Since $f:M\rightarrow M$ is Lipschitz by non-exhaustivity, $f(X)$ is absolutely continuous. Now suppose $p_{f(X)}$ is the probability density function for $f(X)$. Note that by the inverse mapping theorem for Lipschitz maps $D f^{-1}(y)$ exists for almost all $y\in M$. Moreover, the probability density function is by the standard change of variables formula for integrals
\[
p_{f(X)}(x) = p_{X}(f^{-1}(x))\det (D(f^{-1}))(x).
\]
Since $f^{-1}$ is reflection symmetric, and this symmetry is preserved under geodesics, the same will hold for its derivative. Therefore, $p_{f(X)}$ is reflection symmetric.
\end{proof}

It is well known that for any if $M$ is a Riemannian $G$-homogeneous space and if then, action from $G$ acts as isometries, that is it holds that
\begin{equation}
\dist(\Phi p, \Phi q) = \dist(p,q)
\label{eq:liegroupactionlip}
\end{equation}
for all $p,q\in M$ and all $\Phi \in G$. In other words \eqref{eq:liegroupactionlip} shows that the Lie group action is Lipschitz with Lipschitz constant $1$, i.e. non-expansive, and we shall see that (Matrix) Lie group action is reflection symmetric as well.

\begin{lemma}
\label{lem:distancelem}
Let $G$ be a matrix Lie-group with bi-invariant Riemannian metric, and let $M$ be a $G$-homogeneous space with induced Riemannian metric from $G$. Then, for any $p,q\in M$, and any $\Phi \in G$, it holds that if $X = \log_p(q)$, and $\hat{q}=\exp(-X)p$, then
\[
\dist(\Phi p,\Phi q)= \dist(\Phi p, \Phi \hat{q}).
\]
\end{lemma}

\begin{proof}
Suppose $X\in T_eG$ is horizontal such that $q=\exp_{\mathrm{M}}(X)p$ and set $\hat{q}=\exp_{\mathrm{M}}(-X)p$. Then, it holds that
\[
\dist(\Phi q,\Phi p) = \norm{\pi_{\text{Horizontal}}\left( \Phi X \Phi^{-1}\right)}_G
\]
and
\[
\dist(\Phi \hat{q},\Phi p) = \norm{\pi_{\text{Horizontal}}\left( -\Phi X \Phi^{-1}\right)}_G.
\]
Since the horizontal projection is a linear map the result follows.
\end{proof}

Before we go to the main theorem we shall need to check that the space of reflection symmetric variables on Euclidean-space indeed is closed under addition.

\begin{lemma}
\label{lem:convsymm}
Let $X,Y:(\Omega,\mathcal{F},\mathbb{P})\rightarrow \mathbb{R}^n$ be absolutely continuous and independent random variables. If the probability density function of both of them are reflection symmetric around $\mathbb{E}[X]=x$ and $\mathbb{E}[Y]=y$, respectively, then the probability density function of $X+Y$ is reflection symmetric around $x+y$.
\end{lemma}

\begin{proof}
Note that it is classical that the resulting probability density function is
\[
p_{X+Y}(x) = \int_{\mathbb{R}^n} p_X(x-z)p_Y(z) dz
\]
i.e. the convolution of the two prior probability density functions. Also, it holds that
\[
\begin{aligned}
p_{X+Y}(x+y+v) &= \int_{\mathbb{R}^n} p_X(x+y+v-z) p_Y(z) dz&& z\mapsto y-z\\
&= \int_{\mathbb{R}^n} p_X(x+v+z) p_Y(y-z) dz\\
&= \int_{\mathbb{R}^n} p_X(x-v-z) p_Y(y+z) dz&& z\mapsto -z\\
&= \int_{\mathbb{R}^n} p_X(x-v+z) p_Y(y-z) dz&& z\mapsto y-z\\
&= \int_{\mathbb{R}^n} p_X(x+y-v-z) p_Y(z) dz\\
&= p_{X+Y}(x+y-v).
\end{aligned}
\]
\end{proof}

Now we are ready to prove the main theorem.

\begin{theorem}
\label{thm:ARexpect}
Let $M$ be a Riemannian $G$-homogeneous space with $G$ a compact matrix Lie-group and let
\[
Z_\ell = \exp_{Z_{\ell-1}}(\varepsilon_\ell + \log_{Z_{\ell-1}}(\Phi Z_{\ell-1}))
\]
be an AR(1)-process on $M$ as in \eqref{eq:autoregressivehomo}. Suppose that $\varepsilon_\ell$ and $Z_{\ell-1}$ are absolutely continuous w.r.t. $\mathbb{P}$ with respective unique expected values $\mathbb{E}[\varepsilon_\ell]=0$ and $\mathbb{E}[Z_{\ell-1}]$. Moreover, assume that $\varepsilon_\ell$  and $Z_{\ell-1}$ are reflection symmetric around $0\in T_{Z_\ell} M$ and $\mathbb{E}[Z_{\ell-1}]$, respectively, and that they are all mutually independent for all $\ell$. Furthermore, suppose $\rho$ is the injectivity radius of $M$ and assume that $\varepsilon_{\ell}$ is such that $\mathbb{P}(\norm{\varepsilon_\ell} \leq \rho/2)=1$. Then, it holds that $Z_\ell$ is reflection symmetric around $\Phi \mathbb{E}[Z_{\ell-1}]$ and
\[
\mathbb{E}[Z_\ell]= \Phi \mathbb{E} [Z_{\ell-1}].
\]
\end{theorem}

\begin{proof}
Set $x=\mathbb{E}[Z_{\ell-1}]$. First we show that $p_{Z_{\ell}} (z)$ is reflection symmetric around $\Phi x$ from that $\varepsilon_\ell$ and $Z_{\ell-1}$ are reflection symmetric. By Lemma \ref{lem:distancelem} the Lie group action satisfies the assumptions required for $f$ in Lemma \ref{lem:lipreflect} and therefore $p_{\Phi Z_{\ell-1}}$ is reflection symmetric around $\Phi x$. Moreover, the Riemannian logarithm map is a local isometry, and thus we may once again apply Lemma \ref{lem:lipreflect} to $\log_{Z_{\ell-1}}(\Phi Z_{\ell-1})$. Next we have sum of two independent refection symmetric random variables is reflection symmetric by Lemma \ref{lem:convsymm} and hence
\[
\varepsilon_\ell + \log_{Z_{\ell-1}}(\Phi Z_{\ell-1})
\]
is reflection symmetric around $\log_x(\Phi x)$. Lastly $\exp_x$ is a local isometry again and by Lemma \ref{lem:lipreflect} $Z_{\ell-1}$ is reflection symmetric around
\[
\exp_x(\log_x(\Phi x)) =\Phi x =\Phi \mathbb{E}[Z_{\ell-1}]
\]
This follows from that the Riemannian exponential and logarithm are by construction local isometries, Lemma \ref{lem:distancelem} $ii)$ and a similar symmetry argument, this time on $Z_{\ell-1}$, the probability density function of $Z_\ell$ is necessarily centered at $\Phi x$.

Next we need to show that the function $f:M\rightarrow [0,\infty)$ defined by
\[
f(y) = \int_M p_{Z_\ell}(z) \dist(z,y)^2 \dVol_M(z)
\]
is minimal precisely when $y= \Phi x$. To see this, consider the gradient of $f$  w.r.t. $y$. By definition of the distance function, it holds that $\nabla_y\dist(y,z)^2 = 2\log_y$. Therefore,
\[
\nabla f(y)= -\int_M p_{Z_\ell}(z) 2 \log_{y}(z)  \dVol_M(z)
\]
which will be zero at $y= \Phi x$ since $p_{Z_\ell}$ is rotationally symmetric around $\Phi x$ and since $\log_{\Phi x}(z)$ is anti-symmetric around $\Phi x$. Note that
\[
\mathbb{E}[Z_\ell\mid Z_{\ell-1}=u] = \arginf_{m\in M} \int_M \dist(m,z)^2 p_{Z_\ell\mid Z_{\ell-1}=u}(z) \dVol(z)
\]
and by definition it holds that
\[
p_{Z_\ell\mid Z_{\ell -1}=u}(z)=p_{\exp_u(\varepsilon_\ell +\log_{u}(\Phi u))}(z).
\]
Since $\varepsilon_\ell$ is reflection symmetric around $0\in T_u M$ and since $\varepsilon_\ell$ is independent of the event $Z_{\ell-1}=u$, it holds that $p_{Z_\ell\mid Z_{\ell-1}}$ is reflection symmetric around $\Phi u$. This follows from observing that Lemma \ref{lem:convsymm} applies to $\varepsilon_\ell + \log_u(\Phi u)$ and that Proposition \ref{lem:lipreflect} applies to $\exp_u(\varepsilon_\ell + \log_u(\Phi u))$ since $\exp_u$ is a local isometry. Now since $M$ is a Riemannian homogeneous space of a compact Lie group, $M$ is compact and it has bounded sectional curvature, therefore \cite[Theorem 1.2]{karcher1977riemannian} applies using that $\mathbb{P} (\norm{\varepsilon_\ell} \leq \rho/2)=1$ and we can conclude that
\[
\mathbb{E}[Z_{\ell}\mid Z_{\ell-1}=u] =\arginf_{m\in M} \int_{M} \dist(\Phi u,z)^2 p_{Z_\ell\mid Z_{\ell -1}=u}(z) \dVol(z) = \Phi u.
\]

Hence, it holds that
\[
\mathbb{E}[Z_\ell]=\mathbb{E}[\mathbb{E}[Z_{\ell}\mid Z_{\ell-1}=u]] = \mathbb{E}[\Phi Z_{\ell-1}].
\]
Now by a change of variables by the map $h:M\rightarrow M$, defined by $h(u)= \Phi u$, with inverse $h^{-1}(u)=\Phi^{-1} u$, it holds that
\begin{multline*}
\int_M \dist(\Phi v,u)^2 p_{ \Phi Z_{\ell-1}}(u) \dVol(u) \\= \int_M \dist(\Phi v,\Phi u)^2 p_{\Phi Z_{\ell-1}}(\Phi u) \abs{\det(D h^-1)(u)} \dVol(u)\\ = \int_{M} \dist(v, u)^2 p_{Z_{\ell-1}} \dVol(u),
\end{multline*}
where we have used that $\Phi$ acts $M$ as an isometry on $M$ by \eqref{eq:liegroupactionlip} and therefore $\abs{\det(D h^{-1})}=1$. But the last integral is minimal only when $v=x= \mathbb{E}[Z_{\ell-1}]$. Hence,
\[
\mathbb{E}[Z_\ell] = \Phi \mathbb{E}[Z_{\ell-1}].
\]

\end{proof}



\section{Conjugate Gradient descent to solve system identification}
\label{sec:conjgrad}

Here we propose a way of conjugate gradient descent on the group action representation of autoregressive processes on the Orthogonal/Unitary group and some homogeneous spaces over $O(n)$ and $U(n)$ respectively. Note that inherent to this formulation of AR processes, the high degree of multicollinearity of the parameters will make algorithms such as Newton's method impractical since the second derivative of the M.S.E. will be unstable when noise is small.

\subsection{Orthogonal/Unitary groups}
\label{sec:orth}

Here we shall consider the special case when the manifold is the Lie group itself, i.e. $O(n)$. This part is essentially solved by previous works, see \cite{fiori2009algorithm}, but we state it here for completeness. Due to the group structure of the manifold, the parameter of an AR(1) process may immediately be estimated as
\[
\hat{\Phi}_{\ell}=Y_{\ell+1} Y_{\ell}^{-1},
\]
for every $\ell$. Then the average $\hat{\Phi}$ may be computed using a barycenter algorithm on all the $\hat{\Phi}_{\ell}$ as in \cite{pennec2013exponential}. The numerics turn out to work as expected as seen in Figure \ref{fig:Liegroupaverage}.



\subsection{Stiefel Manifolds}
\label{sec:stief}

Consider  $M=\St_{n,k}$, the real compact Stiefel manifold with the natural left $O(n)$-action. We want to minimize the equation $f:O(n)\rightarrow \mathbb{R}_{\geq 0}$ defined by
\[
f(\Phi) = \sum_{\ell=1}^N \dist(Z_{\ell+1},\Phi Z_{\ell})^2.
\]
The issue with solving the minimization problem on Stiefel manifolds for a general $f$ using a Newton algorithm has been treated in \cite{edelman1998geometry}, though this requires a computation of the Euclidean gradient and Hessian of $f$.

For the induced metric (aka canonical metric) on $M$ the geodesic exponential map corresponds to
\[
\exp_X\begin{pmatrix}
A\\ B
\end{pmatrix}= \exp_{\mathrm{M}} \begin{pmatrix}
A & -B^T\\ B& 0
\end{pmatrix} X
\]
where $\exp_{\mathrm{M}}$ denotes the matrix exponential. In contrast to the geodesic exponential map on $O(n)$ the corresponding inverse map is not easy to find. On the other hand we shall make use of the pad\'e approximation of the exponential map to get an extrinsic expression for $f$. Note that
\[
\exp_X(\Delta)=\exp_{\mathrm{M}}(\pi_{Horizontal}\Delta X^T)X \approx \left(I_n-\frac{1}{2}\Delta X^T \right)^{-1} \left(I_n+\frac{1}{2}\Delta X^T\right)X,
\]
see \cite{jauch2020random,yuan2019global,zhu2019matrix}, and thus given $X,Y\in M$ by straightforward algebra it holds that
\[
\log_X(Y)\approx 2(Y-X) (I_k +X^TY)^{-1}.
\]
\begin{remark}
There is a method developed by Ralf Zimmermann, see \cite{zimmermann2017matrix}, which approximates the logarithm map more accurately. The proposed method here could be modified in this regard, though we believe this second order approximation is sufficient for our purposes.
\end{remark}
Since 
\[
\dist(X,Y)^2= \norm{\log_X(Y)}^2 \approx \frac{1}{2} \tr(\log_X(Y)^T \left(I_k-\frac{1}{2} XX^T\right)\log_X(Y))
\]
we have the following approximation of $f$,
\[
\begin{aligned}
f(\Phi) \approx \sum_{j=1}^N \tr &\left(\left(I_k+Y^T_{j}\Phi^T Y_{j+1}\right)^{-1}\left( 3I_k - Y_j^T \Phi^TY_{j+1}-Y_{j+1}^T \Phi Y_j\right.\right.\\ 
& \left. \left.- Y_j^T\Phi^T Y_{j+1} Y_{j+1}^T \Phi Y_j\right) \left( I_k +Y_{j+1}^T \Phi Y_j\right)^{-1} \right).
\end{aligned}
\]

To simplify notation we shall denote
\[
A=\left(I_k + Y_j^T \Phi^T Y_{j+1}\right)^{-1},
\]
\[
B=3I_k -Y_{j+1}^T\Phi Y_j - Y_j^T \Phi^T Y_{j+1} - Y_j^T\Phi Y_{j+1} Y_{j+1}^T \Phi Y_j,
\]
and
\[
C= \left(I_k + Y_{j+1}^T \Phi Y_{j}\right)^{-1}
\]
Consider the inversion map $\inv X\mapsto X^{-1}$ on $\mathrm{GL}(n)$, it holds that
\[
D\inv_X(\Delta) = -X^{-1} \Delta X^{-1}.
\]
The Euclidean derivative of $f$ w.r.t. $\Phi$ is the following linear map
\[
\begin{aligned}
\nabla f_\Phi(\Delta)= -2\sum_{j=1}^N \tr &\left(A  Y_j^T \Delta^T Y_{j+1} AB C+ A \left(  Y_j^T \Delta^T Y_{j+1} +  Y_{j+1}^T \Delta Y_{j}\right.\right.\\
&\left.\left.+Y_j^T\Delta^T Y_{j+1} Y_{j+1}^T \Phi Y_{j} + Y_j^T \Phi^T Y_{j+1} Y_{j+1}^T \Delta Y_{j}\right)C\right.\\
& \left.+ABCY_{j+1}^T \Delta Y_{j}C\right).
\end{aligned}
\]

Thus, we have the following scheme for estimating $\Phi\in O(n)$, using conjugate gradient decent, given $\{Y_j\}_{j=1}^{N+1}\subset \St_{n,k}$ and initial guess $\Phi_0$.

\begin{algorithm}
\caption{Estimating system parameters on $\St_{n,k}$}
\label{alg:stiefelgrad}
\begin{algorithmic}[1]
\State Compute 
\[
\Delta=\sum_{j}^{\dim (O(n))}\nabla f_{\Phi_0}(\Phi_0 e_j) \cdot e_j
\]
where $e_j$ is an ON basis on $\mathfrak{so}(n)$.
\State Find $\tau\in [-1,1]$ such that
\[
f(\Phi_0 \exp(-\tau \Delta))
\]
is minimal.
\State Update $\Phi$ by
\[
\Phi= \Phi_0\exp(-\tau \Delta)
\]
\State Repeat step 1-3 with $\Phi_0=\Phi$ but with $\Delta$ orthogonal to previous $\Delta$ using a Grahm Schmidt procedure, do this a select number of times or until $\text{steps}=\dim(O(n)$.
\State Repeat step 1-4 with $\Phi_0=\Phi$ until tolerance is small enough, e.g. $\norm{\Delta} \abs{\tau}<Tol$.
\end{algorithmic}
\end{algorithm}


Figures \ref{fig:graddescentst},\ref{fig:graddescentnoise},\ref{fig:graddescentstoverd},\ref{fig:graddescentstoverk} show some simulations estimating the system parameter of an AR(1) process on the Stiefel manifolds.

\subsection{Grassmann Manifolds}
\label{sec:grass}

Here we consider the Grassmann manifold $M=\Gr_{n,k}$, we identify $\Gr_{n,k}$ as equivalence classes of $n\times k$ matrices by $\St_{n,k}/O(k)$. Once again given observations $\{Y_j\}_{\ell=1}^{N+1}\subset M$ we want to find $\Phi \in O(n)$ such that
\[
\sum_{\ell=1}^N \dist(Y_{\ell+1},\Phi Y_{\ell})^2 \leq \sum_{\ell=1}^N \dist(Y_{\ell+1},\Psi Y_{\ell})^2 \quad \forall \Psi  \in O(n).
\]
At $[I_{n,k}]$ the tangent space can be identified as
\[
T_{I_{n,k}}\Gr_{n,k}=\left\{\begin{pmatrix}
0\\ V
\end{pmatrix}:V\in \mathcal{M}_{n-k,k}(\mathbb{R})\right\}.
\]
In general it is very difficult to write out an explicit parametrizable tangent vector, so the Question is now; given $\Delta\in T_Y \St_{n,k}$ how do we get $[\Delta]\in T_{[Y]}\Gr_{n,k}$ in order to compute $\dist([X],[Y])$? A general tangent vector in $T_{Y}\St_{n,k}$ may be written as
\[
\Delta = \begin{pmatrix}
A& -B^T\\B& 0
\end{pmatrix} X\qquad A\in \mathfrak{o}(k), B\in \mathcal{M}_{n,k}(\mathbb{R})
\]
and the corresponding Grassmannian tangent (i.e. in horizontal direction) is
\[
[\Delta] = \begin{pmatrix}
0& -B^T\\B& 0
\end{pmatrix} X\qquad B\in \mathcal{M}_{n,k}(\mathbb{R}).
\]
Now since $X^T \Delta=A$ is anti-symmetric we have the following Grassmann metric at the base point $X\in \St_{n,k}$
\[
g_X(V,W)^2=\tr(V^T(I_n-XX^T)W),
\]
note the difference of a factor of $\frac{1}{2}$ compared to the Stiefel case. Moreover, note that this metric is degenerate on the Stiefel manifold but non-degenerate precisely on the Grassmann manifold. Therefore the map we want to minimize can be approximated as follows:
\[
F(\Phi)=\sum_{\ell=1}^N 4 \tr\left(A (I_k-Y_{\ell}^T\Phi^T Y_{\ell+1}Y_{\ell+1}^T \Phi Y_{\ell})C \right)
\]
where $A$ and $C$ are as previously defined. Let
\[
D=I_k-Y_{\ell}^T\Phi^T Y_{\ell+1}Y_{\ell+1}^T \Phi Y_{\ell}.
\] 
Then, the gradient, with respect to $\Phi$, is

\begin{multline*}
\nabla f_\Phi(\Delta)= -\sum_{\ell=1}^N 4 \tr \left(A  Y_\ell^T \Delta^T Y_{\ell+1} AD C+ A \left( Y_\ell^T\Delta^T Y_{\ell+1} Y_{\ell+1}^T \Phi Y_{\ell} \right.\right.\\
\left.\left. + Y_\ell^T \Phi^T Y_{\ell+1} Y_{\ell+1}^T \Delta Y_{\ell}\right)C +ADCY_{\ell+1}^T \Delta Y_{\ell}C\right).
\end{multline*}

Now the algorithm is precisely as done in the case of the Stiefel manifoldl, see Algorithm \ref{alg:stiefelgrad}, but with the above gradient instead. In Figures \ref{fig:graddescentgr},\ref{fig:graddescentgrnoise},\ref{fig:graddescentgroverd},\ref{fig:graddescentgroverk} one can find some simulations estimating the parameters of an AR(1) process using this algorithm.



\appendix
%\fancyhead[LE]{\slshape APPENDIX}



\section{Basic Riemannian Geometry}
\label{sec:basicriemgeom}

To make this paper more accessible, we shall here give a brief introduction to some concepts in differential geometry that are necessary to understand the results of this paper. Everything discussed in this section can be found in \cite{lee2009manifolds} which gives a good modern summary of differential geometry.

Most fundamental is the manifold, its aim is to generalize $\mathbb{R}^n$ to mathematical objects that locally behave like $\mathbb{R}^n$, a framework which allows for analysis (both stochastic and classical) on them. Abstractly, one starts with a set $M$. A \textit{chart} on $M$ is a subset $U_\alpha\subseteq M$ and an injective map $u_{\alpha}:U\rightarrow \mathbb{R}^n$, often denoted as $( U_\alpha,u_\alpha)$. Two charts, $(U_\alpha,u_\alpha)$ and $(U_\beta,v_\beta)$ are said to be (smoothly) compatible either if, $U_\alpha \cap U_\beta = \emptyset$, or if $U_\alpha \cap U_\beta \neq \emptyset$, then the map
\[
u_\alpha \circ (v_\beta\mid_{U_\alpha \cap U_\beta})^{-1} :v_\beta(U_\alpha \cap U_\beta) \rightarrow u_\alpha(U_\alpha \cap U_\beta)
\]
is a diffeomorphism between open subsets in $\mathbb{R}^n$.  An \textit{atlas} is a collection of compatible charts $\mathcal{A}=\{( U_\alpha,u_\alpha)\}_{\alpha \in \mathscr{A}}$, with $\mathscr{A}$ some indexing set, such that
\[
\bigcup_{\alpha \in \mathscr{A}} U_\alpha = M.
\]
Moreover, it is a classical exercise to show that compatibility is an equivalence relation on atlases. The charts allow for differentiation of functions between manifolds. Consider $f:M\rightarrow N$, where $M$ and $N$ are smooth manifolds. Fix $p\in M$. Let $u:U\rightarrow \mathbb{R}^m$ be a chart around $p$, let $v:V\rightarrow \mathbb{R}^n$ be a chart around $f(p)$ and suppose $f(U)\subset V$. Then, $f$ is called \textit{smooth} at $p$ if 
\[
v\circ f \circ u^{-1}:u(U) \rightarrow v(V)
\]
is a smooth function at $u(p)$ in the sense of functions between finite dimensional vector spaces. Next, we have the tangent space.

\begin{definition}
Let $M$ be a smooth manifold and let $p\in M$ be a point. Denote $C_p$ as the set of all curves $\gamma:(-\varepsilon,\varepsilon): \rightarrow M$ such that $\gamma(0)=p$. We say the two curves $\gamma_1, \gamma_2\in C_p$ are equivalent, $\gamma_1 \sim \gamma_2$ if for all smooth $f:M\rightarrow \mathbb{R}$, it holds that
\[
\ddf { }{ }{t} f\circ \gamma_1 \mid_{t=0} = \ddf { }{ }{t} f\circ \gamma_2 \mid_{t=0}.
\]
Then, the \textit{tangent space} at $p$ is the vector space
\[
T_p M := C_p/\sim.
\]
The \textit{tangent bundle} $TM$ is the disjoint union of all the tangent spaces, i.e.
\[
TM = \bigsqcup_{p\in M} T_pM
\]
\end{definition}

Moreover, the \textit{tangent map} of $f:M\rightarrow N$ at $p$, is the map $T_p f:T_pM \rightarrow T_{f(p)}M$ defined by 
\[
\gamma\in C_p \longmapsto f\circ \gamma \in C_{f(p)}.
\]
It is easy to check that this is a well-defined assignment.

In order to solve extremal problems on manifolds, one has the need of (topological) metrics. In differential geometry the topological metric on a manifold is called the \textit{distance}, $\dist$, and it is generated by the so called \textit{Riemannian metric}. For every $p\in M$, the Riemannian metric at $p$ is the linear map
\[
g_p : T_pM \rightarrow T_pM
\]
such that $g_p$ is positive definite and symmetric. Furthermore, as we vary the base point $p$ this assignment of linear maps should be smooth, this means that 
\[
g:p\longmapsto g_p(T_{u(p)}u^{-1} v,T_{\tilde{u}(p)}\tilde{u}^{-1}w), \qquad v,w\in \mathbb{R}^n
\]
is a smooth assignment where $u,\tilde{u}:U\rightarrow \mathbb{R}^n$ are charts around $p$. If $\gamma:[a,b]\rightarrow M$ is a curve on $M$ then, the distance between $p=\gamma(a)$, $q=\gamma(b)$ along $\gamma$ is defined as
\[
\dist_\gamma (p,q)=\int_a^b g_{\gamma(s)}(T_s\gamma,T_s\gamma)^{1/2} ds.
\]
The \textit{geodesic} curve going from $p$ to $q$ is any curve $\tau:[a,b]\rightarrow M$ such that $\dist_\tau(p,q) \leq\dist_\gamma(p,q)$ for all curves $\gamma$ starting at $p$ and ending at $q$. Moreover, the \textit{distance} $\dist(p,q)$ between two points $p,q\in M$ is defined as the distance along the geodesic going from $p$ to $q$. Observe here that the distance function is a topological metric on $M$. Note, that $\tau$ is only guaranteed to be unique when $p$ and $q$ lie in a neighborhood small enough. Since geodesics are locally unique one can define a chart which is locally an isometry between $T_p M$, with metric from the inner product $g_p$ and the distance function on $M$.


\begin{definition}
\label{def:geodexplog}
Let $p\in M$ be a point and let $U\subseteq T_pM$ be open and such that all length minimizing geodesics $[0,\norm{V}]\rightarrow M$ with initial velocity $X= V/\norm{V}$, $V\in U$ are unique. Then, the assignment $\exp_p(V) = \gamma_p(V)$ where $\gamma_p$ is the unique geodesic starting at $p$ with initial velocity $V/\norm{V}$ and length $\norm{V}$, is a map $U\rightarrow M$ called the \textit{Riemannian exponential map} at $p$. It's inverse $\log_p: \exp_p(U)\rightarrow U$ is the \textit{geodesic normal chart} around $p$, and it's associated coordinates are called the \textit{geodesic normal coordinates}. The maximal $\rho_p >0$ such that $\exp_p$ is injective when it is restricted to the ball
\[
B_\rho(0) = \{x\in T_p X: \norm{x}<\rho\}
\]
is referred to as the \textit{injectivity radius of $M$ at $p$}. The minimal injectivity radius at all points on $M$, $\rho =\inf_{p\in M} \rho_p$ is referred to as the \textit{injectivity radius of $M$}.
\end{definition}


The geodesic normal coordinates are of particular importance. In geodesic normal coordinates, the gradient and Hessian are reduced to the Euclidean gradient and Hessian. That is, if $x_i$ are geodesic normal coordinates centered at $p$, i.e. the functions generated by the exponential map and some orthonormal basis on $T_pM$, then 
\[
\nabla f(p) = \frac{\partial f}{\partial x_i}(p) \partial_{x_i}
\]
and
\[
\Delta f(p) = \frac{\partial^2f}{\partial x_i^2} (p).
\]
Utilizing spherical geodesic normal coordinates centered around $p$, 
\[(r,\theta_1,\dots , \theta_{n-1}),\]
i.e. the transformed geodesic normal coordinates $x_i$ pre-composed with the transformation that defines the spherical coordinates in $\mathbb{R}^n$, one has $r=\dist(p,q)$. Moreover, we shall mean by the coordinate $(-r,\theta_1,\dots, \theta_{n-1})$ the reflected vector that lies in the opposite direction of $(r,\theta_1,\dots,\theta_{n-1})$. Moreover, manifolds for which the Riemannian exponential map is defined on all of $T_p M$ is referred to as \textit{geodesically complete}.

A Lie group $G$ is a smooth manifold and a group such that the group multiplication map
\[
\mu:(x,y) \longmapsto x\cdot_G y,
\]
and the inversion map
\[
\inv : x\longmapsto x^{-1}
\]
are smooth maps. The left action $L_p:G\rightarrow G$ is defined by $L_p(q)=\mu(p,q)$ and the right action $R_p:G \rightarrow G$ is defined by $R_p(q)=\mu(q,p)$. Sometimes, there exists a bi-invariant Riemannian metric on $G$, i.e. $g_p$ is such that
\[
g_{p}(v_p,w_p)=g_{\mu(q,p)}(T_p L_q v_p,T_pL_q w_p)=g_{\mu(p,q)}(T_p R_q v_p,T_pR_q w_p).
\]
A very important special case are the matrix Lie groups, which are the Lie groups that are closed subgroups of invertible linear maps $\mathrm{GL}(n)$ over $\mathbb{R}$ or $\mathbb{C}$. In particular, the Orthogonal group $O(n)$ is the real (Matrix) Lie group defined as the set
\[
O(n)=\left\{ A\in \mathrm{GL}(n): AA^{T}= I_n \right\}
\]
together with the multiplication map and inversion map inherited from $\mathrm{GL}(n)$. The tangent space of $O(n)$ at the identity is precisely the Lie algebra $\mathfrak{o}(n)$ which is the space of alternating $n\times n$ matrices. On $O(n)$ we have the invariant bi-invariant metric, also known as the \textit{Frobenious inner product}
\[
g_{e}(X,Y) = \tr (X^T Y),
\]
for $X,Y\in \mathfrak{o}(n)$, then it is extended to all other points through the bi-invariance identity. For this metric, we have that the Riemannian exponential map coincides with the matrix exponential,
\[
\exp_{\mathrm{M}}(X) := \sum_{k=0}^\infty \frac{1}{k!} X^k.
\]
Note that we still distinguish between the two maps by designating the subscript $M$ for the matrix exponential map. 

Finally, one may consider a closed Lie subgroup $H$ of some Lie group $G$ and consider the quotient $G/H$. In general it is not a Lie group, since $H$ may not be normal. But it still has a lot of useful structure from the underlying Lie group. Such quotients are referred to as \textit{homogeneous spaces}. If $G$ is a Lie group with bi-invariant Riemannian metric $g$ one define the \textit{induced metric} on the homogeneous space $G/H$ as restricting the metric to the horizontal part of the tangent space. More precisely since $\mathfrak{h}\subset \mathfrak{g}$ as a vector space we may (orthogonally) decompose any element $X\in \mathfrak{g}$ into horizontal part $\pi(X) \in T_{[e]}G/H$ and a vertical part $V(X)\in \mathfrak{h}$. In particular, by considering $O(n-k)$ as a subgroup of $O(n)$, $k<n$, through the embedding $\iota: O(n-k)\rightarrow O(n)$ defined by
\[
\iota: X \longmapsto \begin{pmatrix}
I_{k} & 0\\ 0 & X
\end{pmatrix}.
\]
One consider the quotient space $\St_{n,k} := O(n)/O(k)$ called the Stiefel manifold and the quotient space $\Gr_{n,k}:=\St_{n,k}/O(k)$ called the Grassmann manifold. Note that here $O(k)$ is considered as the embedded subgroup in $O(n)$ by
\[
X\longmapsto \begin{pmatrix}
X & 0\\ 0 & I_{n-k}.
\end{pmatrix}
\]
Observe that one may interpret the Stiefel manifold as the space of orthogonal $n$ by $k$ matrices and the tangent space of $\St_{n,k}$ at $I_{n,k}$ as the space of $n$ by $k$ matrices of the form
\[
T_{I_{n,k}} \St_{n,k} = \left\{ \begin{pmatrix}
A\\ B
\end{pmatrix} \in \mathcal{M}_{n,k}(\mathbb{R}): A\in \mathfrak{o}(k) \text{ and } B\in \mathcal{M}_{n-k,k} \text{ arbitrary} \right\}.
\]
By considering the induced metric on $\St_{n,k}$ from the underlying metric on $O(n)$ we get a Riemannian structure on $\St_{n,k}$ and the Riemannian exponential map on $\St_{n,k}$ at $I_{n,k}$ is the map
\[
\exp_{I_{n,k}}\left(\begin{pmatrix}
A\\B
\end{pmatrix}\right) = \exp_{\mathrm{M}}\left( \begin{pmatrix}
A &-B^T\\ B & 0
\end{pmatrix}\right)
\]
and the exponential map at some arbitrary point $p\in \St_{n,k}$ is attained by post-multiplication from the right by $p$. Since there are infinitely many representations of $[X]\in \St_{n,k}$ as the equivalence class of matrices $X\in \mathfrak{o}(n)$, the choice of alternating matrix above is referred to as the horizontal projection
\[
\pi_{\text{Horizontal}}: \begin{pmatrix}
A & -B^T\\ B& C
\end{pmatrix} \longmapsto \begin{pmatrix}
A & -B^T\\ B& 0
\end{pmatrix}
\]
where $A\in \mathfrak{o}(k)$, $B\in \mathcal{M}_{n-k,k}$ and $C\in \mathfrak{o}(n-k)$.

\section{Simulations}
\label{sec:simulations}



Here we shall give an outline on some performed simulations for the respective algorithm. All the simulations here are for some special cases of $O(n)$-homogeneous spaces $M$. Moreover, for all cases, we start by  choosing some random element $\Phi \in SO(n)$ with standard deviation of $\dist(\Phi,I_n)\sim 0.01$. The noise terms are generated $N$-times by considering $\varepsilon_j$ as normal distribution on $\mathfrak{o}(n)$ with mean $0\in \mathfrak{o}(n)$ and small standard deviation $\sigma \id_{\mathfrak{o}(n)}$. This may be done by taking the normal distribution on $\mathcal{M}_{n,n}$ and then projecting to the subspace of anti-symmetric matrices. Then starting at some point $p\in M$, the elements $Z_{j}=\exp_{\mathrm{M}}(\varepsilon_j)\Phi Z_{j-1}$ are inductively generated with $Z_{0}=p$. Even though this is not exactly analogous to the formulation in Equation \eqref{eq:liegroupactionAR}, by Baker-Campbell-Hausdorff this is a first order approximation of the latter. Now we apply an algorithm for each case given in Section \ref{sec:conjgrad} which gives an estimate of $\hat{\Phi}$ and then the error is computed by $\dist(\Phi,\hat{\Phi})$.

For the case when $M=O(n)$ we apply the algorithm as outlined in Section \ref{sec:orth}, i.e. through the barycenter algorithm compute the average of all the $Z_{j+1}^{-1} Z_j$. The error between the estimate and $\Phi$ is of the order $\mathcal{O}(\sqrt{N}^{-1})$ with coefficient depending on all other parameters, $n,\Phi,\sigma$. This is expected since this algorithm is a generalization of sample average estimates in Euclidean space, which also has the same order term. In figure \ref{fig:Liegroupaverage} this is illustrated for the case $O(20)$ and $\sigma=0.1$.

\begin{figure}[H]
\centering
\includegraphics[width=0.9\textwidth]{liegroupalgorithmoverN.jpg}
\caption{Simulation an AR(1) process using barycenters on $O_{20}$ where the noise is of approximate amplitude $10^{-1}$.}
\label{fig:Liegroupaverage}
\end{figure}

For the case when $M=\St_{n,k}$ the algorithm given in Section \ref{sec:stief} gives an estimate and the error once again is of the order $\mathcal{O}(\sqrt{N}^{-1})$. This is illustrated in Figure \ref{fig:graddescentst} for the case $\St_{15,5}$.



\begin{figure}[H]
\centering
\includegraphics[width=0.9\textwidth]{graddescentSt.jpg}
\caption{Simulation of the gradient descent method on $\St_{15,5}$ where the noise is of approximate amplitude $10^{-3}$.}
\label{fig:graddescentst}
\end{figure}

In Figure \ref{fig:graddescentnoise} we can see that for $\St_{20,5}$ the error is dependent on the size of the noise, where $N=60$.

\begin{figure}[H]
\centering
\includegraphics[width=0.9\textwidth]{graddescentnoise.jpg}
\caption{Simulation of the gradient descent method on $\St_{20,5}$ over different sizes of the noise.}
\label{fig:graddescentnoise}
\end{figure}

Another observation is how the error term is dependent on the dimensional parameters $n$ and $k$. In Figure \ref{fig:graddescentstoverd} we fix $k=5$ and $N=200$, and we see an increasing tendency of the error when $n=d$ gets larger.

\begin{figure}[H]
\centering
\includegraphics[width=0.9\textwidth]{graddescentStoverd.jpg}
\caption{ Simulation of the gradient descent method on $\St_{d,5}$ over different $d$ where the noise is of approximate amplitude $10^{-3}$.}
\label{fig:graddescentstoverd}
\end{figure}

The opposite behavior can be seen when increasing the parameter $k$ in $\St_{n,k}$ when $k \lesssim n/2$ as seen in Figure \ref{fig:graddescentstoverk} for $\St_{20,k}$ and $N=200$. This is the case even though the dimension of $\St_{n,k}$ increases, the degrees of freedom of the parameter $\Phi$ decreases when $k$ increases.

\begin{figure}[H]
\centering
\includegraphics[width=0.9\textwidth]{graddescentStoverk.jpg}
\caption{Simulation of the gradient descent method on $\St_{20,k}$ over different $k$ where the noise is of approximate amplitude $10^{-3}$.}
\label{fig:graddescentstoverk}
\end{figure}

Lastly, for the case $M=\Gr_{n,k}$, we similarly apply the algorithm in Section \ref{sec:grass} to the simulated data set. Once again the error term is of the order $\mathcal{O}(\sqrt{N}^{-1})$, as seen in the case for $\Gr_{15,5}$ and $N=200$ in Figure \ref{fig:graddescentgr}.

\begin{figure}[H]
\centering
\includegraphics[width=0.9\textwidth]{graddescentGr.jpg}
\caption{Simulation of the gradient descent method on $\Gr_{15,5}$ where the noise is of approximate amplitude $10^{-3}$.}
\label{fig:graddescentgr}
\end{figure}

Furthermore, we have a similar dependency on the noise term $\sigma$, see Figure \ref{fig:graddescentgrnoise} for the case $\Gr_{20,5}$ and $N=60$.

\begin{figure}[H]
\centering
\includegraphics[width=0.9\textwidth]{graddescentGrnoise.jpg}
\caption{Simulation of the gradient descent method on $\Gr_{20,5}$ over different sizes of the noise.}
\label{fig:graddescentgrnoise}
\end{figure}

For the same reason as for the Stiefel manifolds, higher degrees of freedom for $\Phi$ means larger error. Fixing $k$ we see the same behavior as previously, see Figure \ref{fig:graddescentgroverd} for the case $k=5$. 

\begin{figure}[H]
\centering
\includegraphics[width=0.9\textwidth]{graddescentGroverd.jpg}
\caption{Simulation of the gradient descent method on $\Gr_{d,5}$ over different $d$ where the noise is of approximate amplitude $10^{-3}$.}
\label{fig:graddescentgroverd}
\end{figure}

In Figure \ref{fig:graddescentgroverk} we observe the same behavior as in the case of the Stiefel manifold, for the same reason.

\begin{figure}[H]
\centering
\includegraphics[width=0.9\textwidth]{graddescentGroverk.jpg}
\caption{Simulation of the gradient descent method on $\Gr_{20,k}$ over different $k$ where the noise is of approximate amplitude $10^{-3}$.}
\label{fig:graddescentgroverk}
\end{figure}


\printbibliography

\end{document}
\endinput
